% Part: first-order-logic
% Chapter: axiomatic-deduction

\documentclass[../../../include/open-logic-chapter]{subfiles}

\begin{document}

\iftag{FOL}
      {\olchapter{fol}{axd}{Axiomatische \usetoken{P}{derivation}}}
      {\olchapter{pl}{axd}{Axiomatische \usetoken{P}{derivation}}}

\begin{editorial}
  We hebben het materiaal in dit hoofdstuk nog niet aangepast aan de
  verschillende tags. Die tags geven aan welke logische verbindingstekens
  en kwantoren, de tekens voor `voor alle' en `er bestaat', als basistekens
  worden genomen en welke worden gedefinieerd. Voorlopig behandelen we al
  die tekens als basistekens, behalve $\liff$: dat teken behandelen we als
  gedefinieerd. Staat de tag FOL aan, dan maken we een versie met
  kwantoren. Anders maken we er een zonder.
\end{editorial}

\olimport{rules-and-proofs}

\olimport{axioms-rules-propositional}

\iftag{FOL}{%
  \olimport{axioms-rules-quantifiers}
}{}

\olimport{proving-things}

\iftag{FOL}{%
  \olimport{proving-things-quant}
}{}

\olimport{proof-theoretic-notions}

\olimport{deduction-theorem}

\iftag{FOL}{%
  \olimport{deduction-theorem-quantifiers}
}{}

\olimport{provability-consistency}

\olimport{provability-propositional}

\iftag{FOL}{%
  \olimport{provability-quantifiers}
}{}

\olimport{soundness}

\iftag{FOL}{
  \olimport{identity}
}{}

\OLEndChapterHook

\end{document}
